\clearpage
\setcounter{page}{1}
\maketitlesupplementary
\appendix

\section{Timing}
\label{sec:supp_timing}
\begin{table}[h]
  \centering\small
  \begin{tabular}{@{}l rr rr@{}}
\toprule
System & \multicolumn{2}{c}{batch 1} & \multicolumn{2}{c}{batch 64} \\
\cmidrule(lr){2-3}\cmidrule(lr){4-5} & ms & FPS & ms/img & img/s \\
\midrule
DINOv3, 1 view & 14.7 & 68.1 & 2.9 & 347 \\
DINOv3 + orbit & 69.6 & 14.4 & 12.1 & 82 \\
PoE, 1 view & 64.0 & 15.6 & 50.0 & 20 \\
PoE + orbit & 278.5 & 3.6 & 206.8 & 5 \\
\bottomrule
\end{tabular}
  \caption{\textbf{End-to-end latency} of the augmented trees (shared RTX 6000 Ada; upper bounds; image already decoded).}
  \label{tab:timing}
\end{table}
At batch 1 (\cref{tab:timing}), a DINOv3 forward pass takes 6.4\,ms and one tree 7.4\,ms. MoGe-2 adds 21\,ms for the network and 20\,ms for the CPU consensus-mask step per view, which is the main optimisation target. Orbit voting multiplies the cost by $4.4$--$4.7\times$. Batched evaluation reduces tree cost to 0.1--0.3\,ms per image.

\section{Model lineage and fusion}
\label{sec:supp_lineage}
\begin{table}[h]
  \centering\footnotesize
  \setlength{\tabcolsep}{3pt}
  \begin{tabular}{@{}lccc@{}}
    \toprule
    Design (clean trees, PoE greedy) & Synth. & LB & SL \\
    \midrule
    MLP path decoder, hard targets & 6.24 & 31.1 & 43.7 \\
    GRU path decoder & 6.15 & -- & 40.6 \\
    MLP decoder, geodesic soft targets & 6.07 & 28.2 & 42.1 \\
    Transformer decoder, hard targets & 5.69 & 31.6 & 41.8 \\
    \textbf{Transformer + soft targets} & \textbf{5.54} & \textbf{26.9} & \textbf{36.7} \\
    \quad + beam 8 instead of greedy & no gain & -- & -- \\
    \quad + tangent residual at the leaf & $-0.13$ & $\pm0.0$ & $\pm0.0$ \\
    \quad 16$\times$16 instead of 8$\times$8 DINOv3 tokens & 5.56 & 27.0 & 36.0 \\
    \midrule
    DINOv3 tree alone & 5.35 & 37.1 & -- \\
    MoGe-2 tree alone & -- & 43.7 & -- \\
    Learned concat fusion, one Transformer (8$\times$8) & 5.34 & 40.8 & -- \\
    Equal per-level PoE (same grid) & 5.58 & 27.9 & -- \\
    Weighted PoE, $w_{\text{DINO}}{=}0.8$ (proxy-selected) & 5.14 & 32.0 & 36.7 \\
    \bottomrule
  \end{tabular}
  \caption{\textbf{Design steps} on clean (non-augmented) trees, seed means. Dashes: not evaluated in that study.}
  \label{tab:lineage}
\end{table}

\Cref{tab:lineage} lists the design steps that led to the base tree. The Transformer decoder with geodesic soft targets was best; beam search at full data brought no gain, and a tangent residual only refined errors within the correct cell.
Fusion is strongly asymmetric across domains. A learned fusion that feeds concatenated tokens into one Transformer matches DINOv3 alone on synthetic data but is much worse on real images ($40.8^\circ$ vs.\ $27.9^\circ$ for PoE on the same grid). We interpret this as the learned fusion over-fitting render statistics, while training-free PoE keeps the branches' errors partly independent.
A weighted PoE selected on the rotated validation proxy ($w{=}0.8$) is $5^\circ$ worse on lightbox than equal weights. Following our protocol, we did not use real data to choose another weight.
The consensus foreground mask reduces MoGe-2's background leakage on real images from 0.52 to 0.12 (lightbox) and from 0.57 to 0.12 (sunlamp). It is what makes the MoGe-2 branch useful.

\section{Negative results}
\label{sec:supp_negative}
\begin{table}[h]
  \centering\footnotesize
  \setlength{\tabcolsep}{2.5pt}
  \begin{tabular}{@{}lccc@{}}
    \toprule
    Approach (on clean trees unless noted) & Synth. & LB & SL \\
    \midrule
    Clean PoE tree (reference) & 5.54 & 26.9 & 36.7 \\
    $SO(3)$ flow posterior on tree memory, KDE mode & 5.22 & 29.8 & 37.6 \\
    \quad best of 32 samples (\emph{oracle}) & 2.60 & 9.7 & 10.6 \\
    Flow re-noise verifier over beams + flips + corners & 4.78 & 26.5 & 32.7 \\
    \quad best candidate (\emph{oracle}) & 2.25 & 7.9 & 11.4 \\
    Kurz corner tree~\cite{kurz2017tesseract} & 5.20 & 28.8 & 39.1 \\
    Tree distilled from flow masses ($\beta{=}1$, 1 seed) & 5.62 & 26.2 & 37.7 \\
    Discrete flow matching on tree paths & 7.22 & 29.8 & 39.8 \\
    \midrule
    Aug.\ DINOv3 tree (reference) & 4.84 & 22.6 & 22.1 \\
    \quad + $SO(3)$ flow, KDE & 4.10 & 25.2 & 24.4 \\
    \bottomrule
  \end{tabular}
  \caption{\textbf{Refiners and verifiers that did not fix real flips} (seed means). Oracles are not deployable and show headroom only.}
  \label{tab:negative}
\end{table}

\Cref{tab:negative} collects approaches that improved synthetic accuracy but not, or only slightly, real accuracy.
A Riemannian flow-matching posterior~\cite{chen2024rfm} conditioned on the tree's encoder memory contains good samples: its best-of-32 oracle reaches $9.7^\circ$ on lightbox. It also assigns the ground truth a much higher log-likelihood than the tree does (3.85 vs.\ 0.98 nats on lightbox), and its 90\% regions over-cover (0.98). Yet its mode estimate is worse than the tree's.
Re-scoring an expanded candidate set (beam, $180^\circ$ flips, cell corners) by re-noise/return consistency helps synthetic data and sunlamp, but does not repair lightbox flips. Synthetic render-and-compare templates also failed to choose the right flip.
A discrete flow on the tesseract path~\cite{gat2024dfm} was clearly worse. We note that interpolating path tokens is not geometry-preserving: 45--53\% of intermediate paths lie more than $30^\circ$ from both endpoints.
Distilling flow masses into the tree did not pass its criterion.
An oracle that stops decoding at the best depth would gain only $0.51^\circ$ on synthetic data, so adaptive depth was not pursued.
These results support the diagnosis of \cref{sec:diagnosis}: the right mode is reachable, but scorers trained on clean synthetic features cannot select it on real images. Lighting-robust training and multi-view consistency can.

\section{Flip transitions on SPEED+}
\begin{table}[h]
  \centering\small
  \setlength{\tabcolsep}{2.6pt}
  \begin{tabular}{@{}lcccc@{}}
\toprule
System & \multicolumn{2}{c}{Lightbox} & \multicolumn{2}{c}{Sunlamp} \\ \cmidrule(lr){2-3}\cmidrule(lr){4-5}
 & fixed & broken & fixed & broken \\
\midrule
Aug.\ tree, 1 view & 51\% of 401 & 1.0\% & 59\% of 206 & 1.0\% \\
Aug.\ tree + orbit & 62\% of 401 & 0.4\% & 80\% of 206 & 0.4\% \\
Aug.\ DINOv3 + orbit & 52\% of 401 & 1.3\% & 74\% of 206 & 0.5\% \\
\bottomrule
\end{tabular}
  \caption{\textbf{Flip transitions relative to the clean PoE tree.} \emph{Fixed}: share of the clean tree's ${>}150^\circ$ errors that become ${<}20^\circ$ (denominator given). \emph{Broken}: share of its ${<}20^\circ$ predictions that become ${>}150^\circ$.}
  \label{tab:flips}
\end{table}

\section{Orbit distillation and 6-DoF details}
\begin{table}[h]
  \centering\footnotesize
  \setlength{\tabcolsep}{2.4pt}
  \begin{tabular}{@{}lll ccc ccc@{}}
\toprule
Arm & Feat. & $\hat t$ & \multicolumn{3}{c}{$E_T$ (m)} & \multicolumn{3}{c}{$E_R$ ($^\circ$), orbit} \\
\cmidrule(lr){4-6}\cmidrule(lr){7-9} & & & Syn. & LB & SL & Syn. & LB & SL \\
\midrule
Box only & -- & -- & 0.561 & 0.547 & 0.553 & -- & -- & -- \\
\midrule
Frozen & DINOv3 & 1 view & 0.090 & 0.233 & 0.240 & 4.21 & 17.40 & 15.27 \\
 &  & orbit & 0.068 & 0.205 & 0.205 &  &  &  \\
Frozen & PoE & 1 view & 0.108 & 0.206 & 0.232 & 3.87 & 12.30 & 12.25 \\
 &  & orbit & 0.089 & 0.184 & 0.206 &  &  &  \\
\midrule
Joint & DINOv3 & 1 view & 0.078 & 0.209 & 0.212 & 4.19 & 17.31 & 14.87 \\
 &  & orbit & \textbf{0.059} & 0.185 & \textbf{0.183} &  &  &  \\
Joint & PoE & 1 view & 0.111 & 0.196 & 0.221 & 3.86 & 12.11 & 11.73 \\
 &  & orbit & 0.093 & \textbf{0.175} & 0.194 &  &  &  \\
\bottomrule
\end{tabular}
  \caption{\textbf{SPEED+ 6-DoF} with the SPACE-HOP translation target, ground-truth crops. Rotation uses PoE or DINOv3 orbit voting. Translation ($\hat t$) is predicted from one view or as the orbit average; seed std of $E_T$ ${\le}\,0.005$\,m. \emph{Box only} uses no image: the object centre is placed at the box centre and the depth at the training mean of $U_z$.}
  \label{tab:sixdof}
\end{table}

\begin{table}[h]
  \centering\small
  \setlength{\tabcolsep}{2.5pt}
  \resizebox{\linewidth}{!}{\begin{tabular}{@{}l cc ccc@{}}
\toprule
Training objective & \multicolumn{2}{c}{1 view, DINOv3} & \multicolumn{3}{c}{PoE + orbit} \\
\cmidrule(lr){2-3}\cmidrule(lr){4-6} & Syn. & Syn.-rot. & Syn. & LB & SL \\
\midrule
Base recipe (GT loss) & 4.84\,{\scriptsize$\pm$0.05} & 9.88\,{\scriptsize$\pm$0.22} & 3.87 & 12.30 & 12.25 \\
+ paired views & 4.96\,{\scriptsize$\pm$0.06} & 10.04\,{\scriptsize$\pm$0.21} & 3.94 & 12.38 & 12.45 \\
+ leaf-level transported consistency & 4.59\,{\scriptsize$\pm$0.04} & 9.43\,{\scriptsize$\pm$0.14} & 3.71 & 12.41 & 11.74 \\
+ hierarchical transported consistency & 4.65\,{\scriptsize$\pm$0.03} & 9.57\,{\scriptsize$\pm$0.06} & 3.74 & 12.49 & 11.73 \\
\midrule
+ orbit distillation & 4.57\,{\scriptsize$\pm$0.05} & 9.33\,{\scriptsize$\pm$0.18} & \textbf{3.69} & \textbf{11.21} & \textbf{11.09} \\
\bottomrule
\end{tabular}}
  \caption{\textbf{Training-time consistency.} Single-view DINOv3 trees on synthetic and rotated synthetic validation (left) and PoE + orbit (right). The consistency losses transport each view's beam distribution into another view's frame via \cref{eq:roll}.}
  \label{tab:consistency}
\end{table}

Retention of the orbit gain by the single-view students, defined as (base greedy $-$ student greedy)/(base greedy $-$ base orbit), is 43\%/18\%/22\% (DINOv3) and 30\%/28\%/23\% (PoE) on synthetic/lightbox/sunlamp, and 13--15\% on the rotated proxy.
For the joint 6-DoF model, the ESA-style score (rotation error in radians plus normalised translation error, with the challenge precision floors) is 0.084/0.241/0.238 for PoE.

\section{Per-object results}
\begin{table}[h]
  \centering\footnotesize
  \begin{tabular}{@{}l cc@{}}
\toprule
Object & PoE+orbit & DINOv3+orbit \\
\midrule
\multicolumn{3}{@{}l}{\emph{T-LESS}} \\
\quad 1 & 23.82 & 25.00 \\
\quad 4 & 14.50 & 15.69 \\
\quad 5 & 6.22 & 7.07 \\
\quad 11 & 11.90 & 12.92 \\
\quad 18 & 14.38 & 14.78 \\
\quad 21 & 34.78 & 37.24 \\
\quad 27 & 10.36 & 12.43 \\
\multicolumn{3}{@{}l}{\emph{SYMSOL-T}} \\
\quad cone & 1.83 & 1.79 \\
\quad cube & 2.72 & 2.90 \\
\quad cyl & 1.67 & 1.65 \\
\quad icosa & 3.81 & 3.71 \\
\quad tet & 3.35 & 3.06 \\
\multicolumn{3}{@{}l}{\emph{YCB-V}} \\
\quad 003\_cracker\_box & 85.27 & 74.61 \\
\quad 006\_mustard\_bottle & 97.96 & 87.29 \\
\quad 011\_banana & 61.63 & 83.23 \\
\quad 025\_mug & 24.25 & 22.76 \\
\quad 035\_power\_drill & 41.34 & 25.07 \\
\quad 037\_scissors & 169.07 & 123.72 \\
\bottomrule
\end{tabular}
  \caption{\textbf{Per-object test error} of the single-label (A0) trees (3-seed mean; symmetry-aware for T-LESS and SYMSOL-T; YCB-V trained on real videos only).}
  \label{tab:perobj}
\end{table}
On T-LESS (\cref{tab:perobj}), most of the test error comes from object 21 (no symmetry metadata, near-symmetric; 18\% of predictions above $90^\circ$) and object 1 (continuous symmetry). Continuous symmetries are evaluated on a $10^\circ$-step discretisation of the symmetry set, which adds up to $5^\circ$ of artificial error for objects 1 and 4 and for the cone and cylinder.

\section{YCB-V diagnosis}
\label{sec:supp_ycbv}
On YCB-V we checked four things:
\begin{enumerate}
\item The trees fit their training frames: $3.8^\circ$ and $3.9^\circ$ on mustard bottle and mug.
\item Test poses are covered by training: the median nearest training rotation is about $5^\circ$ away for cracker box, mustard bottle and banana, whose test errors are still $75$--$87^\circ$.
\item Every crop contains all of the object's label pixels.
\item The ground-truth rotations agree with the segmentation: 65--98\% of the projected model points fall on the object mask, against 28--83\% for a transposed-rotation control.
\end{enumerate}
The failure is therefore memorisation of the 16--20 cluttered training videos per object. The object covers only 5--40\% of the crop and the DINOv3 tokens are not masked. This is not a label or pipeline error. Standard YCB-V training adds the synthetic \texttt{data\_syn} renders, whose backgrounds are black. We pasted each render onto a real frame drawn only from the training sequences (never validation or test), capped at 12,000 crops per object as for the real data, and retrained with the unchanged recipe. \Cref{tab:ycbv_perobj} shows the effect per object.
\begin{table}[h]
  \centering\footnotesize
  \begin{tabular}{@{}l cc cc@{}}
\toprule
 & \multicolumn{2}{c}{real videos} & \multicolumn{2}{c}{+ synthetic} \\
\cmidrule(lr){2-3}\cmidrule(lr){4-5} Object & DINOv3 & PoE & DINOv3 & PoE \\
\midrule
003\_cracker\_box & 77.7 & 78.7 & 5.1 & 5.7 \\
006\_mustard\_bottle & 101.5 & 102.4 & 6.1 & 7.6 \\
011\_banana & 79.9 & 77.6 & 6.1 & 5.3 \\
025\_mug & 24.6 & 28.4 & 6.5 & 5.6 \\
035\_power\_drill & 30.3 & 38.2 & 4.3 & 4.4 \\
037\_scissors & 113.2 & 161.5 & 49.4 & 36.9 \\
\bottomrule
\end{tabular}
  \caption{\textbf{YCB-V per-object test error} ($^\circ$, single view, 3 seeds) without and with synthetic training data.}
  \label{tab:ycbv_perobj}
\end{table}

\section{Zero-shot transfer of the frozen SPEED+ trees}
\label{sec:supp_zeroshot}
\begin{table}[h]
  \centering\footnotesize
  \setlength{\tabcolsep}{3pt}
  \begin{tabular}{@{}l cc ccc@{}}
\toprule
Dataset & \multicolumn{2}{c}{Raw $E_R$} & \multicolumn{3}{c}{Frame-free pair error} \\
\cmidrule(lr){2-3}\cmidrule(lr){4-6} & frozen & chance & frozen & chance & constant \\
\midrule
T-LESS & 117.5 & 103.6 & 95.4\,{\scriptsize$\pm$3.0} & 104.5 & 88.7 \\
SYMSOL-T & 49.4 & 53.7 & 47.8\,{\scriptsize$\pm$0.5} & 53.7 & 53.8 \\
YCB-V & 115.6 & 125.8 & 101.6\,{\scriptsize$\pm$2.8} & 125.8 & 57.8 \\
\bottomrule
\end{tabular}
  \caption{\textbf{Zero-shot}: frozen SPEED+ PoE trees with orbit voting, no retraining, test split. The pair error $\min_S\angle(\hat R_i\hat R_j^\top, R_iSR_j^\top)$ over random same-object pairs cancels the unknown offset between object frames. \emph{Chance}: uniformly random predictions. \emph{Constant}: one rotation predicted for every image.}
  \label{tab:zeroshot}
\end{table}
Object frames are arbitrary, so raw errors are expected near chance. The frame-free pair error (\cref{tab:zeroshot}) is therefore the meaningful measure. On T-LESS and YCB-V it is \emph{worse} than predicting one constant rotation, because both datasets have strong pose priors. On SYMSOL-T, where rotations are uniform, the trees beat chance only because of one shape: the cube ($27.8^\circ$ vs.\ $41.2^\circ$), plausibly because the spacecraft body is box-like. There is no general zero-shot transfer: the trees must be retrained per object.

\section{Disclosures and reproducibility}
\begin{itemize}
\item \textbf{Late feature-equivalence check.} The augmented and rotated views are extracted by a compact feature extractor. A pre-registered check against the full feature caches was run after the main results. It failed its strict element-wise tolerance: the maximum standardised difference was 0.51 for MoGe-2 features (p99 0.004; correlation 0.999998). Its functional follow-up showed that no PoE prediction changed on 512 crops, and that single-branch greedy leaves changed for 1/512 (DINOv3) and 3/512 (MoGe-2) crops.
\item \textbf{SYMSOL-T bug.} Nine MoGe-2 SYMSOL-T trees first trained to NaN: always-empty border cells had a standard deviation near $10^{-6}$, and float16 overflowed. They were retrained with a standard-deviation floor of $10^{-2}$; all other trees were unaffected. Only the corrected evaluation is reported.
\item \textbf{YCB-V paths.} A file-path error in the YCB-V manifest (colour images are JPEG) was fixed before any YCB-V result.
\item \textbf{Post-hoc control.} The constant-prediction control in \cref{tab:zeroshot} was added after the T-LESS zero-shot result and is labelled as post hoc.
\end{itemize}
Every number in the results tables is generated by a script directly from the evaluation logs of the named runs.
